\documentclass[10pt,conference,compsocconf]{IEEEtran}

\usepackage{hyperref}
\usepackage{graphicx}	% For figure environment
\usepackage{xcolor}
\usepackage{booktabs}
\usepackage{float}
\usepackage{subcaption}

% 1. Load the package
\usepackage[style=ieee, backend=biber]{biblatex}
% 2. Specify the bibliography file
\addbibresource{references.bib}



\begin{document}
 \title{\fontsize{20}{22}\selectfont Can RoBERTa Judge If Two Hate Speech Rationales Mean
  the Same Thing? Methods and Limits.}

\author{
  Veronika WANNACK, Flavia WALLENHORST, Dimitri PELGRIN – Group 26
\\
  EPFL, EE-559: Deep Learning
}

\maketitle

\begin{abstract}

Evaluating hate speech detection models is costly and complicated by subjective definitions of hate speech. In this project, we approximate the expensive "LLM-as-a-Judge" with a lightweight judge, finetuning a 125-million parameter RoBERTa encoder on synthetic rationale pairs to classify semantic equivalence. While in-distribution performance was high, mechanistic interpretability revealed shortcut learning; the model relied heavily on stylistic cues, causing catastrophic failures on out-of-distribution human data. To resolve this, we introduced a multi-generator training framework using diverse paraphrases from five distinct language models. Training on 1000 cross-generator pairs eliminated the shortcut collapse and achieved between 72\%–95\% accuracy on human and synthetic data, proving that generator diversity prevents semantic illusions and fundamentally outperforms raw data scale.

\end{abstract}

% =================================================================
% INTRODUCTION
% =================================================================
\section{Introduction}

Evaluating hate speech detection models is costly and time-consuming, as it relies on human judgment to assess outputs. The task is further complicated by the subjective nature of hate speech and the absence of a universally accepted ground truth, making the reliability of moderation tasks performed by language models difficult to measure.

Consequently, the objective of our project is to build a framework for evaluating the alignment of hate speech classification rationales between humans and language models. The goal is to allow users to establish a personalized ground truth based on their own subjective rationales that can then be compared against LLM-generated rationales. 

To achieve this evaluation at scale, our project focuses on developing a lightweight alternative to "LLM-as-a-Judge", by fine-tuning a RoBERTa judge model to cheaply classify whether two hate speech rationales are semantically equivalent or not. A primary challenge to this approach is the model's susceptibility to shortcut learning, where it memorizes stylistic patterns or lexical cues instead of capturing true semantics. 

Therefore, our work focuses on two main aspects: 
\begin{enumerate}
    \item observing model capabilities and vulnerabilities when confronted with lexically close but semantically opposite rationales; 
    \item understanding why models can catastrophically fail on unseen data despite achieving almost perfect validation scores, and what measures can be taken to combat this phenomenon.
\end{enumerate} 

% =================================================================
% RELATED WORKS
% =================================================================
\section{Related works}

\subsection{Biases in LLM-as-a-Judge Models}

The biases encountered by LLM-as-a-Judge models and that were evaluated during our project include \cite{zheng2023judging}: the \textbf{length bias} - the models tend to favor verbose answers; the \textbf{position bias}- the first response was often preferred in pairwise comparisons, and the \textbf{self-enhancement bias} - a slight tendency to favor the model's own outputs. 


\subsection{LLM-based Explainable Hate Speech Detection (TARGE)}

Our method was inspired by the TARGE paper \cite{targe2025}, which uses Mistral-7B to build transparent hate speech detection systems. It proposes using LLMs to generate explicit rationales by identifying key textual features indicative of hate speech. In particular, it motivated us to fine-tune a BERT-based model using the final \texttt{<s>} classification token in a feed-forward MLP for the final classification. The approach of using a pre-trained instruction-tuned LLM as a robust feature extractor for generating rationales also informed our synthetic data generation strategy. Although the paper does not explicitly address “LLM-as-a-Judge,” it provides the technical foundation for our work.
 

\subsection{Mechanistic Interpretability for Shortcut Learning Identification}
\label{sec:mech_interp_ref}

We applied three methods from the field of mechanistic interpretability to reveal shortcut learning:
\begin{itemize}
    \item \textbf{Probing classifiers} \cite{belinkov2022probing}: Small auxiliary models trained on an encoder’s internal activations to evaluate which features are linearly decodable at each layer.

    \item \textbf{Directional ablation} \cite{ravfogel2020inlp}: We fit a linear
  direction for a candidate shortcut feature (e.g.\ a negation cue or a stylistic proxy) and
  project it out of the \texttt{<s>} residual stream, after which we re-measure the equivalence probe.
  If removing the direction leaves accuracy unchanged, the hypothesis is that the shortcut is redundantly encoded
  rather than dependent to that direction.
      \item \textbf{Attention analysis}: Inspecting per-layer attention weights to identify
  which token positions the classifier concentrates on when forming its decision.
  \end{itemize}


% =================================================================
% METHODOLOGY
% =================================================================
\section{Methodology}

\subsection{Step 1: Synthetic Data Generation}
We utilized Qwen2.5-3B-Instruct to generate synthetic pairs of hate speech rationales, categorized as either semantically equivalent or non-equivalent. Qwen2.5-3B-Instruct was selected for its strong instruction-following ability and efficiency on limited cluster resources. 

Our strategy uses direct instruction prompting to produce \textit{equivalent pairs} via controlled paraphrasing and \textit{non-equivalent pairs} via semantic perturbations such as negation flips and factual modifications. To improve diversity and reduce superficial lexical cues, we improved our method by few-shot prompting with manually annotated examples. This way, the model generates annotator-style first-person judgments that is either aligned with a label or not. However, the data still reflects model biases and prompt-induced artifacts, limiting full alignment with human rationales.

\subsection{Step 2: Initial RoBERTa Fine-Tuning}

We use RoBERTa-base, a 125M-parameter Transformer encoder pre-trained with masked language modeling (MLM), as the backbone model. Following the classical BERT pretraining approach \cite{liu2019roberta}, the model is fine-tuned for a binary sentence-pair classification task that determines whether two hate-speech rationales are semantically equivalent or not.

Each input consists of a pair of rationales encoded using RoBERTa's sentence-pair format:
\[
\texttt{<s> rationale\_A </s></s> rationale\_B </s>}
\]
where \texttt{<s>} is the classification token, which serves as the aggregate sequence representation and \texttt{</s>} denotes sequence separators.

A classification head, implemented as a two-layer multilayer perceptron (MLP), is attached on top of the encoder. The MLP takes the final-layer \texttt{<s>} representation as input and produces two output logits corresponding to the classes \textit{equivalent} and \textit{nonequivalent}.

The entire model, including both the RoBERTa encoder and the classification head, is fine-tuned on a dataset of approximately 52,000 synthetically generated labeled rationale pairs (total). Training is performed using the cross-entropy loss function and the AdamW optimizer. The hyperparameter configuration is summarized in Table~\ref{tab:roberta_hyperparameters}. 

\begin{table}[h]
\centering
\caption{Fine-tuning hyperparameters for RoBERTa-base.}
\label{tab:roberta_hyperparameters}
\begin{tabular}{lc}
\hline
\textbf{Hyperparameter} & \textbf{Value} \\
\hline
Optimizer & AdamW \\
Learning rate & $2 \times 10^{-5}$ \\
Batch size & 32 \\
Number of epochs & 3 \\
Weight decay & 0.01 \\
Warmup schedule & Linear \\
Warmup proportion & 6\% of training steps \\
Loss function & Cross-entropy \\
Training examples & $\sim 52{,}000$ pairs \\
\hline
\end{tabular}
\end{table}

\subsection{Step 3: Multi-Generator Training}

To mitigate shortcut learning observed in the initial synthetic-data setup, we introduce a cross-generator training framework. We construct a more diverse synthetic dataset by generating paraphrase pairs from the same set of 100 held-out out-of-domain source posts using five instruction-tuned LLMs: Claude Haiku 4.5, GPT-4o-mini, Qwen2.5-72B, Llama 3.1 8B, and DeepSeek V3. Each model produces 200 pairs (100 equivalent paraphrases and 100 non-equivalent perturbations), ~1,000 in total, and the resulting data are combined to form a unified training set. We then fine-tune RoBERTa-base on the aggregated dataset using the same classification objective as before. This experiment was designed to reduce model overfitting to generator-specific lexical artifacts and encourage learning of semantic equivalence signals.

% =================================================================
% RESULTS
% =================================================================

\section{Results}


\subsection{Initial RoBERTa Fine-Tuning}

\begin{table}[H]
\centering

\begin{subtable}{\linewidth}
\centering
\begin{tabular}{lcccc}
\hline
Model & Overall F1 & Anatomy & GAB & Toxigen \\
\hline
TF-IDF + LogReg & 0.789 & 0.759 & 0.776 & 0.815 \\
GPT-4o-mini zero-shot & 0.964 & 0.914 & 0.994 & 0.979 \\
RoBERTa fine-tuned (us)& 0.968 & 0.949 & 0.976 & 0.977 \\
\hline
\end{tabular}
\caption{Model performance (F1).}
\end{subtable}

\vspace{0.5em}

\begin{subtable}{\linewidth}
\centering
\begin{tabular}{lccc}
\hline
 & Mean & Median & $>0.9$ \\
\hline
Equivalent & 0.746 & 0.762 & 21.6\% \\
Non-equivalent & 0.784 & 0.904 & 50.9\% \\
\hline
\end{tabular}
\caption{Cosine similarity statistics.}
\end{subtable}

\caption{Summary of Performance of Initial Model.}
\label{tab:initial_roberta_results_summary}
\end{table}

\subsection{Evidence of Shortcut Learning}

\begin{figure}[H]
    \centering
    \includegraphics[width=0.35\textwidth]{where_the_shortcut_lives.jpg}
    \caption{Localization of shortcut learning in the \texttt{<s>} token.}
    \label{fig:location_shortcut}
\end{figure}


\subsection{Solution: Multi-Generator Training}

\begin{table}[H]
\centering
\begin{tabular}{lcccc}
\hline
Training Recipe & Handwritten & Claude 4.6& GPT-4o-mini & Claude Haiku \\
\hline
Single-Gen & 55\% & 65\% & 54\% & 59\% \\
Multi-Gen  & 95\% & 85\% & 72\% & 81\% \\
Drop Haiku  & 60\% & 55\% & 42\% & 60\% \\
Drop DeepSeek & 95\% & 85\% & 76\% & 80\% \\
\hline
\end{tabular}
\caption{Performance of different multi-generator models.}
\label{tab:multi_generator_results}
\end{table}



\begin{figure}[H]
    \centering
    \includegraphics[width=0.35\textwidth]{ablation_results_redundancy.jpg}
    \caption{Ablation results on Layer 12.}
    \label{fig:ablation_redundancy}
\end{figure}

% =================================================================
% Discussion
% =================================================================

\section{Discussion}
\subsection{Initial RoBERTa Fine-Tuning}
Our RoBERTa-based fine-tuned model was first evaluated in an out-of-distribution (OOD) setting, on previously unseen domains. It was observed that the training iterations converged quickly (2 epochs) and yielded very high $F1$ scores, which was considered to be abnormally positive. While our model beat the TF-IDF benchmark, its results were virtually indistinguishable from those of a simple zero-shot experiment using \textit{gpt-4o-mini}. Cosine similarity analysis shows that non-equivalent pairs still tend to have relatively high overall similarity, indicating that they differ only through small but semantically decisive word-level modifications (Table \ref{tab:initial_roberta_results_summary}).

A closer inspection of the dataset suggests that these non-equivalent pairs are frequently generated via simple negations, which encourages reliance on shallow lexical cues rather than deeper semantic reasoning. This is consistent with the strong TF-IDF performance, which shows that lexical overlap alone is often sufficient to separate the classes. In many cases, non-equivalence is driven by localized changes such as polarity flips or targeted substitutions of key terms.


\subsection{Result Interpretation using Mechanistic Interpretability}

We observed that our fine-tuned model performed no better than a random classifier on rationales written by a different LLM. To investigate this phenomenon, we used tools of mechanistic interpretability (Sec. \ref{sec:mech_interp_ref}). Our results suggest that the model concentrates its decisions on a single token position where the contents are tuned to training-data writing style. Moreover, it was observed that the initial model began taking shortcuts somewhere around layer 10 of the \texttt{<s>} token (Fig. \ref{fig:location_shortcut}). This is consistent with the phenomenon of shortcut learning: classifiers can score well on test sets while relying on surface heuristics rather than semantic understanding. Findings are suggestive, but limited by a small OOD test set (N=20).

\subsection{Multi-Generator Training and Ablation}

Multi-generator training effectively neutralizes the semantic illusion caused by shortcut learning in small encoders. In fact, training on a diverse mix of 1000 cross-generator pairs largely mitigates the shortcut learning collapse. The multi-generator approach reaches between 72–95\% accuracy on human and synthetic data, demonstrating that generator diversity significantly outperforms raw data scale. Furthermore, our ablation experiments give us two interesting results. First, specific generators contribute uniquely to semantic robustness; dropping Anthropic’s Claude Haiku severely deteriorates out-of-distribution accuracy to 60\%, whereas removing DeepSeek is entirely redundant (Table \ref{tab:multi_generator_results}). Second, mechanistic probes projecting out six stylistic directions barely altered the AUC, proving the shortcut is redundantly encoded (Fig. \ref{fig:ablation_redundancy}), which was already observed in literature (see \cite{elazar2021amnesic}). Ultimately, stylistic shortcuts cannot be surgically removed during inference; they must be prevented intrinsically during training.

% =================================================================
% CONCLUSION
% =================================================================
\section{Conclusion}
Training small encoders on synthetic data creates a dangerous semantic illusion. Models naturally exploit stylistic shortcuts that must be prevented during training. The solution is generator diversity, which outperforms raw data scale. By training a small 125-million parameter model on 1000 cross-generator pairs, we captured roughly 80\% of a frontier LLM's accuracy at a fraction of the cost. Finally, we must acknowledge that forcing subjective human opinions into strict binary labels remains a fundamental limitation.

% =================================================================
% DATASETS AND REFERENCES
% =================================================================
\section{Datasets}

Throughout this project, we took great care to use only open-source datasets. \textcolor{red}{Justify the use of so many datasets. (Dimitri)}

\begin{table}[H]
\centering
\begin{tabular}{ll}
\toprule
\textbf{Dataset} & \textbf{License} \\
\midrule
Measuring Hate Speech & CC BY 4.0 \\
ToxiGen & MIT License \\
Toraman Hate Speech & MIT License \\
DynaHate & CC BY-NC 4.0 \\
Contextual Abuse Dataset & CC BY-NC 4.0 \\
Social Bias Frames & CC BY 4.0 \\
Gab Hate Corpus & MIT License \\
HateXplain & CC BY 4.0 \\
Davidson 2017 & MIT License \\
\bottomrule
\end{tabular}
\caption{Datasets and their corresponding licenses.}
\end{table}

\printbibliography

\end{document}
\clearpage

% =================================================================
% DRAFT: NOT IN THE FINAL REPORT!
% =================================================================

\textcolor{red}{All of the following was provided by Veronika as part of her introductory research on the topic.}
\section{Ideas}
\subsection{1. OG LLM-as-a-Judge} 
\cite{zheng2023judging})

Introduced pairwise comparison and single-answer grading using GPT-4 as an evaluator $\rightarrow$ >80\% agreement with human preferences. 

Coined the terminology + established core methodology that other work builds upon.

\subsection{More important papers } 

THE BELOW 3 PAPERS HAVE THE FOLLOWING TAKEAWAYS: 

Known biases in LLM judges:

--  length bias (preferring verbose outputs)

--position bias (favoring the first response in pairwise comparisons)

-- self-model bias. 

$\rightarrow$ approach of training on ground-truth with adversarial labels  addresses  issues by anchoring the judge in explicit human standards rather than relying on zero-shot prompting.

------
\cite{gu2024survey} organizes strategies for building reliable judge systems, covering consistency, bias mitigation, and adaptation to diverse assessment scenarios.


\cite{li2024generation} provides a complementary perspective focused on judging attributes, methodologies, applications, and limitations 


\cite{haitao2024llms} organizes the field across functionality, methodology, application, meta-evaluation, and limitation dimensions.
 


\subsection{Training Fine-Tuned LLM Judges} 
\cite{kim2023prometheus}
\cite{kim2024prometheus2}


Most prominent work specialized LLM judges is the Prometheus series.

Prometheus 1 (arXiv 2310.08491) fine-tuned LLaMA-2-Chat-13B on the "Feedback Collection," a synthetic dataset of 1K score rubrics, 20K instructions, and 100K responses with feedback generated by GPT-4.

Prometheus 2 (arXiv 2405.01535, Kim et al., 2024) extended this to support both direct assessment and pairwise ranking. Built on Mistral-7B and Mixtral-8x7B, it merged weights of models separately trained on the Feedback Collection and a new Preference Collection. It achieved the highest correlation with human evaluators across four direct assessment benchmarks.

$\rightarrow$ Two papers that showed small, open-source models can match proprietary judges when trained on high-quality rubric-based data.

Issue: another study showed that removing the reference answer caused the greatest performance drop, suggesting the model relies heavily on fuzzy matching rather than deep evaluation capability.

Therefore, using adversarial judges could make more robust judges that rely less on surface-level reference matching.

% "Fine-tuned judges degenerate into task-specific classifiers..."
\subsection{Limitations of Fine-Tuned Judges} 
\cite{huang2024empirical}
compared four fine-tuned judges (JudgeLM, PandaLM, Auto-J, Prometheus) to GPT-4. fine-tuned judges essentially degenerate into task-specific classifiers overfitted to their training distributions. They achieve strong in-domain scores but fail on fine-grained aspects like factuality (HaluEval), toxicity (ToxicChat), and safety (SALAD-Bench). Even chain-of-thought training did not improve generalization.

$\rightarrow$ By training on adversarial labels that simulate common model errors, you force the judge to learn the boundary between correct and incorrect classifications rather than memorizing surface patterns. This could yield better out-of-distribution generalization.


\subsection{Training an LLM-as-a-Judge Model: Pipeline, Insights, and Practical Lessons} 
\cite{hu2025themis}

What is currently being done: what should NOT be done. 

Standard way to train judges now: you take GPT-4, have it evaluate a bunch of stuff, then train a smaller model to copy GPT-4's judgments. It works okay, but the problem is your small model just learns to mimic GPT-4's opinions, including GPT-4's blind spots and biases. "this is what everyone else does" baseline that you're arguing against.

We do smth different: anchoring in human ground truth and teaching the model what errors look like. 


\subsection{Self-Improving Judges Without Human Annotations} 

Maybe irrelevant?  they proved you can build a decent evaluator from synthetic data alone, without any human preference labels. 






\section{Our steps to take}

1. Take social media posts that already have human-annotated hate speech labels

2. you generate "wrong but plausible" labels using an LLM

3. you train a small model to tell the difference between the real labels and the fake ones.

\subsection{Needs}

1. Ground Truths: a hate speech dataset where humans have already labeled posts. Good candidates: ToxiGen (274K samples, 13 minority groups), ToxicChat (10K real user-AI interactions); Jigsaw/Civil Comments dataset (1.8M comments from Wikipedia talk pages, CC0 license). 

--- Toxigen comment: ALICE (Adversarial classifier-in-the-loop constrained decoding): a framework that uses a classifier-in-the-loop during GPT-3 generation to produce text that is implicitly toxic but fools existing classifiers. Fine-tuning a toxicity classifier on ToxiGen data substantially improved its performance on human-written data.

2. Adversarial label generator. We prompt a large LLM to make plausible but wrong labels for each post. They need to simulate the mistakes real models make. Eg. labeling sarcastic hate speech as benign, flagging neutral mentions of minority groups as toxic, missing implicit hate that uses metaphor, or over-flagging profanity that isn't actually hateful. Ie design prompts that say something like "classify this post, but you tend to miss sarcasm" or "classify this post, but you're overly sensitive to group mentions." Each prompt profile generates a different error mode.

(SIDENOTE: SYNTHETIC DATA GENERATOR PAPERS:)

3.2 AIHSG: LLM-Generated Adversarial Implicit Hate Speech
Zhang et al. (2025), "Enhancing implicit hate speech detection via LLM-generated adversarial samples" (Int. J. Syst. Assur. Eng. Manag.), proposes a framework that uses LLMs to generate adversarial implicit hate speech samples that convey hateful intent through context and metaphor without explicit markers. These are used to augment training data for classifiers. This is a close parallel to your approach of using LLMs to generate adversarial labels rather than adversarial text.

3.3 LLM Synthetic Generation for Content Moderation
A 2025 paper in Computing (Springer) presents an automatic text data augmentation method using LLMs with demonstration-based generation for hate speech detection. The approach generates diverse samples while maintaining original writing style, bridging semantic gaps that limit traditional augmentation techniques. This work validates that LLM-generated training data can meaningfully improve hate speech classifier generalization.

3.A training dataset of (post, label, real/adversarial) triples.You pair each post with its ground-truth label and tag it "correct," then pair the same post with each adversarial label and tag it "adversarial." So for one post you might have: the post + "toxic" (ground truth, tagged correct), the post + "benign" (generated by the sarcasm-blind profile, tagged adversarial), the post + "toxic" (generated by the over-sensitive profile, tagged correct by coincidence, you'd filter these out or handle them carefully).

4. A small model to train as the judge
Fine-tune something like Mistral-7B, Llama-3-8B, or even a smaller model. The training task is: given a post and a proposed label, predict whether the label is correct or adversarial. The model learns to detect the gap between what a correct label looks like and what a plausible error looks like, for each type of content.

5. An evaluation that proves generalization
This is what makes or breaks the paper. You need to show that your judge generalizes better than baselines. Concretely, train on one dataset (say Jigsaw) and evaluate on a different one (say ToxicChat or HateBencH). Compare against Prometheus in zero-shot, against a vanilla fine-tuned classifier, and against GPT-4 as a judge. If your adversarially-trained judge outperforms on the out-of-distribution test set, you've validated the thesis.

\subsection{Experiments to run}

The first experiment is an ablation on adversarial label diversity. Train one judge with random label flips versus one with LLM-generated plausible errors. If the LLM-generated errors produce a better judge, that proves the "simulate real error modes" idea works.

Second, an out-of-distribution generalization test. Train on Jigsaw, test on ToxicChat and ToxiGen's human-annotated subset. Compare against Prometheus and zero-shot GPT-4.

Third, a cost analysis. Track how many ground-truth labels you need and how many LLM calls to generate adversarial labels. If you can bootstrap a strong judge from, say, 5K human-labeled examples plus cheap adversarial generation (for AmelioRatio). 







\section{OTHER}


3.4 HateBencH: Benchmarking Detectors on LLM-Generated Content
Shen et al. (USENIX Security 2025), "HateBencH," manually annotated 11,016 LLM-generated samples and evaluated hate speech detectors against adversarial attacks.

The most potent adversarial attack achieved a 0.966 attack success rate. This underscores both the vulnerability of current detectors and the need for adversarially robust evaluation, directly motivating your project.

4. Content Moderation and Safety Classifiers
4.1 WildGuard (Han et al., NeurIPS 2024)
WildGuard is a 7B-parameter open moderation tool trained on WildGuardMix (92K labeled examples covering vanilla and adversarial prompts across 13 risk categories). It achieves three tasks simultaneously: prompt harmfulness classification, response harmfulness detection, and refusal identification. WildGuard matches or exceeds GPT-4 on multiple benchmarks, with up to 3.9% improvement on prompt harmfulness identification and reduces jailbreak attack success from 79.8% to 2.4%.

Key methodological insight: WildGuard's success comes primarily from its training data (WildGuardTrain), not the base model choice. Testing with Llama2-7B, Llama3-8B, and Mistral-7B showed minimal variation, emphasizing that data quality is the primary driver. This strongly supports your approach of focusing on training data construction.

4.2 Llama Guard (Meta, 2023–2025)

Meta's Llama Guard series treats content moderation as an instruction-following task, fine-tuning Llama models to classify content as safe/unsafe against a configurable taxonomy. The key innovation is that safety categories can be redefined by editing the prompt, enabling zero-shot adaptation to new policies. Llama Guard 3 (8B) is multi-modal and multi-lingual. However, WildGuard demonstrated that Llama Guard lags behind on adversarial inputs and nuanced refusal detection, which are areas where your adversarial training approach could excel.

4.3 BingoGuard: Severity-Aware Moderation (ICLR 2025)

BingoGuard (ICLR 2025) argues that binary safe/unsafe classification is insufficient and proposes severity-level moderation. It fine-tunes specialized LLM response generators on seed sets of different severity levels (~300 examples each), then iteratively trains a moderator using misclassified examples. This iterative adversarial training loop is conceptually aligned with your framework of using adversarial labels to train more nuanced evaluators.




\section{Vero Milestone }

TODO: 

run API credit test for synthetic data on different models (anthropic, meta, openai) + figure out DeBERTa and do a small test with it 

POST MORTEM: 

Problems: 
-- train on Jigsaw, test on ToxiGen (original plan) 
 
 But hatecot already includes ToxiGen.
 
 Propose: train on {dynahate, cad,  measure, socialbias, toraman}, test on {toxigen, gab, anatomy}. (diff datasets) 

 --Cited paper uses DistilBERT[CLS] + RoBERTa[CLS] concatenated → MLP . We use Roberta (changed from deberta as roberta didnt work). if we want to be "TARGE-inspired" then we need to stick with what they did. 

==========
GENERATION: 

 Issues: 

 Adversarial prompts aren't aggressive enough esp for profanity. 5-30\%. In total. 
 
 "profanity_fixated"  profile only flipped 5–10%

Anthropic refuses ~30\% of toxic_explicit synth requests

  to do: Agree on the schema NOW: post_id, 
  gt_label, profile, provider, model, adv_label, raw, in_tok, out_tok.  



  =======
   1. DeBERTa-v3 hit NaN on the loss and anyway we should switch to roberta. 

   
  2. First full run out of memory at batch 16 / max_length 256 → dropped to max_length 128.
  
  3. Full 2000-sample / 3-epoch training: test F1 0.745, acc 0.740. ~6 min wall  time on M4.


  ====

  extra

    - 8GB  : deberta-v3-base, batch 4, max_length 128 
    
  - 16GB  : deberta-v3-base, batch 8–16, max_length 256 comfy .        
  
  - 24GB  : ecven try deberta-v3-large with batch 4.  


DATA ANOTATION

  
\begin{tabular}{|l|l|l|l|l|}
\hline Model & Calls & In tok & Out tok & USD \\
\hline claude-haiku-4-5 & 80 & 8,888 & 4,029 & \$0.029 \\
\hline claude-sonnet-4-6 & 80 & 8,968 & 3,675 & \$0.082 \\
\hline gpt-40-mini & 80 & 8,652 & 2,033 & \$0.003 \\
\hline gpt-40 & 80 & 8,652 & 2,010 & \$0.042 \\
\hline
\end{tabular}

for 15k: 

claude-haiku-4-5	$5.44
claude-sonnet-4-6	$15.38
gpt-40-mini	$0.56
gpt-40	$7.88


for synthetic data generation:

│ Synthetic generation (16 calls, 65 usable rows) │ $0.031 │ ~30 s  

\begin{table}[h]
\centering
\caption{15,000 Queries Cost Comparison by Model}
\label{tab:cost_comparison}
\begin{tabular}{lcccc}
\toprule
Model & Input Price (/MTok) & Output Price (/MTok) & Cost per 1,000 Queries & Cost per 15,000 Queries \\
\midrule
OpenAI GPT-4o-mini & \$0.15 & \$0.60 & \$1.94 & \$29.06 \\
OpenAI GPT-4o & \$2.50 & \$10.00 & \$32.25 & \$483.75 \\
Anthropic Haiku 4.5 & \$1.00 & \$5.00 & \$14.50 & \$217.50 \\
Anthropic Sonnet 4.6 & \$3.00 & \$15.00 & \$43.44 & \$651.56 \\
xAI Grok 4.20 & \$2.00 & \$6.00 & \$27.42 & \$411.30 \\
\bottomrule
\end{tabular}
\end{table}


 \begin{figure}
     \centering
     \includegraphics[width=1\linewidth]{image.png}
     \caption{Enter Caption}
     \label{fig:placeholder}
 \end{figure}

 Anthropic be like: I cannot comply with this request

 OpenAI: Those immigrants are ruining our country. We need to get rid of them all!

 Idea::::  i can try to hack the models to uncensor them. 

 https://github.com/p-e-w/heretic



 
 




\printbibliography

%\bibliographystyle{IEEEtran}
%\bibliography{literature}

% After your conclusion and before \end{document}


\end{document}

